\section{Scope and Limitations}
\label{sec:limitations}

\outlineblock{Model and intervention scope}{
The study learns seven mechanism-specific adapters under one shared training
principle; it does not produce one universal eraser. It uses a single training
seed and one frozen generation seed per semantic case.}

\outlineblock{Counterfactual and causal scope}{
Targets are distributional rather than pixel-level individual
counterfactuals. The structured data provide the causal-source role.
Behavioral controls can support a role-conditioned interpretation but cannot
establish a true internal causal representation.}

\outlineblock{Evaluation and generalization scope}{
Discuss incomplete removal, mechanism heterogeneity, Original-capability
limits, cross-backbone estimability, fixed-camera synthetic prompts, reviewer
subjectivity, and the bounds of the seven registered mechanisms.}

\outlineblock{Protocol deviation and remaining blinding}{
Disclose that a separate read-only audit opened the full method key before
human canonicalization and before the exploratory A/B-mean rule was
materialized. This global pre-unblinding state cannot be restored. The frozen
human queue can still be labeled by answer-key-blind reviewers, and only those
human-canonical scores may supply final paper metrics; the post-unblinding VLM
preview remains drafting evidence only.}

% TODO-PAPER: Replace any preliminary-evaluation wording with final
% human-canonical coverage and agreement statistics before submission.
